% !PRINTED 169
\weberpage{169}

%<<CONT
if one sets $\alpha=\beta=\frac12$, $\gamma=1$. One of its
particular integrals is therefore
\weq{10}{F(\tfrac12,\tfrac12,1,x) = 1 + \sum_{\nu=1}^{\infty} \Bigl(\frac{1.3.5\ldots 2\nu-1}{2.4\ldots 2\nu}\Bigr)^2 x^\nu}
\[
{}= \sum_{\nu=0}^{\infty} \frac{\Pi(2\nu)^2}{\Pi(\nu)^4}\Bigl(\frac{x}{16}\Bigr)^\nu,
\]
and this series is convergent as long as the absolute value of
$x$ is less than 1. As the second particular integral one may
choose
\weq{11}{F(\tfrac12,\tfrac12,1,1-x),}
which, however, has a different domain of convergence. If we think of
$x$ as a complex variable represented in a plane, then (10) converges
in a circle described about the origin with radius 1, and (11) in an
equal circle described about the point 1, so that the common domain of
convergence consists of a lune containing the part of the real
axis that lies between 0 and 1. But one can also, and this is important
for us, set up the second particular integral in a form which converges
in the same circle as the series (10). Here it is to be observed
that the second particular integral becomes infinite for $x=0$.

This second particular integral is obtained by the
following passage to the limit. If we denote the hypergeometric
series after \gesperrt{Gauss} by
\[
F(\alpha,\beta,\gamma,x) = 1 + \frac{\alpha.\beta}{1.\gamma}x + \frac{\alpha(\alpha+1)\beta(\beta+1)}{1.2.\gamma(\gamma+1)}x^2 + \cdots
\]
\[
{}= \frac{\Pi(\gamma-1)}{\Pi(\alpha-1)\Pi(\beta-1)} \sum_{\nu=0}^{\infty} \frac{\Pi(\alpha+\nu-1)\Pi(\beta+\nu-1)}{\Pi(\gamma+\nu-1)\Pi(\nu)} x^\nu,
\]
then in general one obtains the two particular integrals
of (9):
\[
\begin{aligned}
y_1 &= F_1(x) = F(\alpha,\beta,\gamma,x), \\
y_2 &= F_2(x) = x^{1-\gamma}(1-x)^{\gamma-\alpha-\beta} F(1-\alpha,1-\beta,2-\gamma,x).
\end{aligned}
\]
But since for $\alpha=\beta=\frac12$, $\gamma=1$ these become identical
with one another, one first sets $\alpha=\beta=\frac12$, $\gamma=1+\varepsilon$, hence
\[
\begin{aligned}
y_1 &= F(\tfrac12,\tfrac12,1+\varepsilon,x), \\
y_2 &= \Bigl(\frac{1-x}{x}\Bigr)^\varepsilon F(\tfrac12,\tfrac12,1-\varepsilon,x).
\end{aligned}
\]
